\documentclass[a4paper,12pt]{article}
\usepackage{amsmath,amssymb}
\usepackage[russian]{babel}
\usepackage[utf8]{inputenc}
\usepackage[T2A]{fontenc}
\usepackage{geometry}
\geometry{margin=1.5cm}

\begin{document}

\section*{Контрольные вопросы и ответы}

\subsection*{Вопрос 1: В каких случаях применяется численное интегрирование?}

Численное интегрирование применяется в следующих случаях:
\begin{enumerate}
    \item Когда подынтегральная функция задана таблично (не известна аналитическая формула).
    \item Когда функция имеет сложное аналитическое выражение и интеграл не берется в элементарных функциях.
    \item Когда первообразная функции не выражается через элементарные функции (например, $e^{-x^2}$, $\sin(x)/x$).
    \item Когда требуется высокая точность вычисления интеграла с контролем погрешности.
\end{enumerate}

\subsection*{Вопрос 2: На чем основано численное интегрирование?}

Численное интегрирование основано на замене интеграла конечной суммой (квадратурной суммой):
\[
\int_a^b f(x)\,dx \approx \sum_{i=0}^{n} c_i f(x_i),
\]
где $x_i$ --- узлы интегрирования, $c_i$ --- весовые коэффициенты. Точность приближения зависит от количества узлов и выбора метода.

\subsection*{Вопрос 3: Что такое квадратурные формулы?}

Квадратурные формулы --- это формулы для приближенного вычисления определенного интеграла в виде взвешенной суммы значений функции в конечном числе точек. Основные примеры: формулы прямоугольников, трапеций, Симпсона, Ньютона-Котеса, Гаусса.

\subsection*{Вопрос 4: Каким образом связана задача численного интегрирования и интерполяция?}

Численное интегрирование часто основано на интерполяции. Подынтегральная функция заменяется интерполяционным полиномом (обычно полиномом Лагранжа), который проходит через узлы интегрирования. Затем интеграл вычисляется от полинома вместо исходной функции. Чем выше степень полинома, тем выше точность аппроксимации.

\subsection*{Вопрос 5: Как оценивается погрешность квадратурной формулы?}

Погрешность квадратурной формулы оценивается несколькими способами:
\begin{enumerate}
    \item \textbf{По порядку аппроксимации:} $R \sim O(h^p)$, где $h$ --- шаг разбиения, $p$ --- порядок.
    \item \textbf{По правилу Рунге:} $R = \frac{|I_{2n} - I_n|}{2^k - 1}$, где $k$ --- порядок метода.
    \item \textbf{Аналитически:} используются формулы для оценки остатка (например, для трапеций: $R = -\frac{(b-a)^3}{12n^2}f''(\xi)$).
\end{enumerate}

\subsection*{Вопрос 6: Какие частные случаи формулы Ньютона-Котеса Вы знаете?}

Частные случаи формулы Ньютона-Котеса:
\begin{enumerate}
    \item $n=1$: формула трапеций
    \item $n=2$: формула Симпсона (парабол)
    \item $n=3$: формула 3/8 Симпсона
    \item $n=6$: используется в данной работе с коэффициентами $[41, 216, 27, 272, 27, 216, 41]$
\end{enumerate}

\subsection*{Вопрос 7: Как называется метод численного интегрирования, в котором подынтегральная функция заменяется полиномом нулевой степени?}

Это методы прямоугольников (левых, правых и средних). Полином нулевой степени --- это константа. В левых прямоугольниках берется значение в левом конце подотрезка, в правых --- в правом, в средних --- в середине подотрезка.

\subsection*{Вопрос 8: В каком методе численного интегрирования подынтегральная функция заменяется квадратичным полиномом?}

Это метод Симпсона (метод парабол). Функция на каждой паре соседних подотрезков заменяется параболой (полиномом 2-й степени), проходящей через три узла. Это дает более высокую точность ($O(h^4)$) по сравнению с методом трапеций ($O(h^2)$).

\subsection*{Вопрос 9: Чем отличается метод трапеций от метода Симпсона?}

\begin{enumerate}
    \item \textbf{Аппроксимирующая функция:} трапеции используют линейные функции, Симпсон --- квадратичные (параболы).
    \item \textbf{Порядок точности:} трапеции имеют $O(h^2)$, Симпсон --- $O(h^4)$.
    \item \textbf{Формула трапеций:} $I_T = \frac{h}{2}[f(x_0) + 2\sum_{i=1}^{n-1}f(x_i) + f(x_n)]$
    \item \textbf{Формула Симпсона:} $I_S = \frac{h}{3}[f(x_0) + 4\sum f(x_{нечетные}) + 2\sum f(x_{четные}) + f(x_n)]$
    \item \textbf{Требование:} для Симпсона $n$ должно быть четным.
\end{enumerate}

\subsection*{Вопрос 10: В чем суть метода средних прямоугольников?}

Метод средних прямоугольников состоит в следующем:
\begin{enumerate}
    \item Интервал $[a,b]$ разбивается на $n$ равных подотрезков длины $h = (b-a)/n$.
    \item На каждом подотрезке $[x_i, x_{i+1}]$ вычисляется значение функции в середине: $f(x_{i+1/2})$.
    \item Интеграл приближается суммой площадей прямоугольников: $I_M = h\sum_{i=1}^{n} f(x_{i-1/2})$.
    \item Погрешность метода: $O(h^2)$.
\end{enumerate}

\subsection*{Вопрос 11: Что представляет собой формула для вычисления интеграла методом трапеций?}

Формула трапеций имеет вид:
\[
\int_a^b f(x)\,dx \approx \frac{h}{2}[f(x_0) + 2f(x_1) + 2f(x_2) + \cdots + 2f(x_{n-1}) + f(x_n)],
\]
где $h = (b-a)/n$, $x_i = a + ih$. Первое и последнее значения берутся с коэффициентом 1, остальные --- с коэффициентом 2. Это соответствует замене функции ломаной линией и вычислению площади под ней как суммы площадей трапеций.

\subsection*{Вопрос 12: Как определить число разбиений интервала интегрирования, используя неравенство для оценки абсолютной погрешности для метода трапеций?}

Погрешность метода трапеций оценивается как:
\[
|R| \leq \frac{(b-a)^3}{12n^2} \max_{x \in [a,b]} |f''(x)|.
\]

Требуя $|R| \leq \varepsilon$, получаем условие на $n$:
\[
n \geq \sqrt{\frac{(b-a)^3 \max |f''(x)|}{12\varepsilon}}.
\]

Таким образом, число разбиений должно быть не менее чем корень квадратный из правой части неравенства.

\subsection*{Вопрос 13: Что такое правило Рунге?}

Правило Рунге --- это метод апостериорной оценки погрешности численного интегрирования. Вычисляются приближения интеграла при числе разбиений $n$ и $2n$, а затем погрешность оценивается как:
\[
R = \frac{|I_{2n} - I_n|}{2^k - 1},
\]
где $k$ --- порядок метода ($k=1$ для прямоугольников, $k=2$ для трапеций, $k=4$ для Симпсона). Правило Рунге позволяет автоматически выбирать число разбиений для достижения требуемой точности.

\subsection*{Вопрос 14: Каким образом можно уменьшить погрешность решения при численном интегрировании?}

Погрешность можно уменьшить следующими способами:
\begin{enumerate}
    \item \textbf{Увеличить число разбиений $n$:} погрешность уменьшается как $O(h^p) = O((1/n)^p)$.
    \item \textbf{Выбрать более точный метод:} метод Симпсона ($O(h^4)$) точнее, чем трапеции ($O(h^2)$).
    \item \textbf{Использовать адаптивные методы:} увеличивать разбиение в местах, где функция быстро меняется.
    \item \textbf{Применять правило Рунге:} автоматически выбирать $n$ до достижения требуемой точности.
    \item \textbf{Использовать методы высокого порядка:} формулы Ньютона-Котеса, Гаусса и т.д.
\end{enumerate}

\subsection*{Вопрос 15: Когда удобнее пользоваться квадратурной формулой Гаусса?}

Квадратурная формула Гаусса удобнее всего использовать в следующих случаях:
\begin{enumerate}
    \item Когда требуется очень высокая точность при умеренном числе узлов.
    \item Когда функция гладкая и хорошо приближается полиномом высокой степени.
    \item Когда можно выбирать расположение узлов (они не обязаны быть равноотстоящими, в отличие от Ньютона-Котеса).
    \item Когда первообразная сложна для вычисления, но функция легко вычисляется в произвольных точках.
\end{enumerate}

Формула Гаусса дает точность $O(h^{2n})$ с использованием $n+1$ узлов, что значительно лучше, чем формулы Ньютона-Котеса того же порядка.

\end{document}
